
Enfin, on remarque que @@M0@@, donc que les morphismes @@M1@@ et @@M2@@ sont obtenus par passage au quotient à partir des morphismes @@M3@@ et @@M4@@ de 3.7. On a donc @@M5@@.

D'après VI 5.8, @@M6@@ pour tout @@M7@@, puisque @@M8@@, @@M9@@ étant de rang @@M10@@. On en déduit donc un homomorphisme :
@@M11@@
et un homomorphisme analogue après tensorisation par @@M12@@. L'homomorphisme @@M13@@ est donc obtenu à partir de @@M14@@ (3.7) par passage au gradué associé. Comme @@M15@@ est inverse à gauche de @@M16@@, @@M17@@ est inverse à gauche de @@M18@@, ce qui prouve en même temps l'injectivité de @@M19@@.

Pour montrer que @@M20@@, il faut montrer que si on a @@M21@@, et @@M22@@, tels que :
@@M23@@
alors il existe @@M24@@ tel que :
@@M25@@

Posons :
@@M26@@
Alors d'après VI 5.8, @@M27@@. On remarque ensuite que @@M28@@ conserve la filtration (en utilisant des résultats de l'exposé VIII, ce qui est sans inconvénient, car 4.8 n'est pas utilisé pour prouver ce fait). En effet, il résulte de 1.8 ii) que @@M29@@ est un morphisme d'intersection complète de dimension virtuelle relative nulle (VIII 4.2). Comme @@M30@@ est projectif, le résultat est conséquence de VIII 3.2. En appliquant @@M31@@ à la relation
